\section{Results}
\label{sec:results}

It is important to note that a total of 2,418 patients met the inclusion criteria (Figure~\ref{fig:flow}). Baseline characteristics are presented in Table~\ref{tab:baseline} --- the groups were well balanced, highlighting the robustness of the propensity-score matching approach \cite{austin2011introduction}.

The primary outcome occurred in 142 patients (11.7\%) in the intervention group and 198 patients (16.4\%) in the control group (hazard ratio [HR] 0.69; 95\% CI 0.56--0.86; $p = 0.001$). Notably, this effect was remarkably consistent across all prespecified subgroups, underscoring the crucial role of early intervention (Figure~\ref{fig:forest}).

The adjusted model was specified as
\begin{equation}
  \log h_i(t) = \log h_0(t) + \beta_1 \text{Treatment}_i + \boldsymbol{\gamma}^\top \mathbf{x}_i ,
  \label{eq:cox}
\end{equation}
where $\mathbf{x}_i$ denotes the vector of baseline covariates listed in Table~\ref{tab:baseline} \cite{cox1972regression, therneau2000modeling}.

Furthermore, secondary outcomes showed a similar pattern, with all-cause mortality of 7.9\% versus 8.6\% (HR 0.92; 95\% CI 0.70--1.21) -- a finding that serves as a testament to the overall safety of the strategy.
